% ============================================================
% SKELETON: Classical Thermodynamics of Simple Fluids
% ------------------------------------------------------------
% Planned position: a new part between "Statistical Derivation of
% Ensembles" (Sections 1-5) and "Phase Transitions" (Sections 6-7).
% Not yet included in main.tex. To preview, add % ============================================================
% SKELETON: Classical Thermodynamics of Simple Fluids
% ------------------------------------------------------------
% Planned position: a new part between "Statistical Derivation of
% Ensembles" (Sections 1-5) and "Phase Transitions" (Sections 6-7).
% Not yet included in main.tex. To preview, add % ============================================================
% SKELETON: Classical Thermodynamics of Simple Fluids
% ------------------------------------------------------------
% Planned position: a new part between "Statistical Derivation of
% Ensembles" (Sections 1-5) and "Phase Transitions" (Sections 6-7).
% Not yet included in main.tex. To preview, add % ============================================================
% SKELETON: Classical Thermodynamics of Simple Fluids
% ------------------------------------------------------------
% Planned position: a new part between "Statistical Derivation of
% Ensembles" (Sections 1-5) and "Phase Transitions" (Sections 6-7).
% Not yet included in main.tex. To preview, add \include{thermo.tex}
% under \part{Classical Thermodynamics} and move that part in front of
% \part{Phase Transitions}.
%
% Conventions of this skeleton:
%   - Each \paragraph has one or two sentences that state the goal of the
%     paragraph, followed by the boxed target result where one exists.
%   - "% TODO" comments list the derivation steps still to be written.
%   - "% FROM" comments name the earlier result that a step reuses, so the
%     chapter does not repeat Sections 1-5.
% ============================================================

\section{Classical Thermodynamics of Simple Fluids}

\paragraph{Scope of this section.}
Sections~3--5 derive the laws of thermodynamics, the potentials and the
fluctuation formulas from entropy maximization. This section turns these
results into measurable quantities and applies them to model fluids. The
response functions and the stability conditions of this section are the
tools of the phase-transition sections, and the ideal gas and the van der
Waals gas are the first explicit models of the notes.

% ============================================================
\subsection{Response Functions}
% ============================================================

\paragraph{Measurable second derivatives.}
Experiments measure how one variable responds to a change of another,
not the potentials themselves. Each \textcolor{blue}{response function} is a
second derivative of a potential.
% FROM: C_V and Var(U) = k_B T^2 C_V (Section 3); delta Q = dH at fixed P
%       (Section 4); G(T, P, N) and its derivatives (Section 5).
% TODO: define C_P = T (dS/dT)_P = (dH/dT)_P, alpha_P, kappa_T,
%       kappa_S = -(1/V)(dV/dP)_S; write each as a second derivative of G.
\begin{equation}
    \boxed{C_P=-T\frac{\partial^2G}{\partial T^2},\qquad
    \alpha_P=\frac{1}{V}\frac{\partial^2G}{\partial T\,\partial P},\qquad
    \kappa_T=-\frac{1}{V}\frac{\partial^2G}{\partial P^2}.}
\end{equation}

\paragraph{Maxwell relations.}
The mixed second derivatives of a potential do not depend on the order of
differentiation. Each potential therefore gives one
\textcolor{blue}{Maxwell relation} between two response-type derivatives.
% TODO: derive from dU, dF, dH, dG (Sections 3-5); note which relation turns
%       an unmeasurable entropy derivative into a measurable one.
\begin{equation}
    \boxed{
    \begin{aligned}
        \left(\frac{\partial T}{\partial V}\right)_{\!S,N}&=-\left(\frac{\partial P}{\partial S}\right)_{\!V,N}, &
        \left(\frac{\partial S}{\partial V}\right)_{\!T,N}&=\left(\frac{\partial P}{\partial T}\right)_{\!V,N},\\
        \left(\frac{\partial T}{\partial P}\right)_{\!S,N}&=\left(\frac{\partial V}{\partial S}\right)_{\!P,N}, &
        \left(\frac{\partial S}{\partial P}\right)_{\!T,N}&=-\left(\frac{\partial V}{\partial T}\right)_{\!P,N}.
    \end{aligned}}
\end{equation}

\paragraph{The cyclic relation.}
Three variables linked by one equation of state satisfy the cyclic
relation. The slope of an isochore in the $P$--$T$ plane then follows from
$\alpha_P$ and $\kappa_T$.
% FROM: the implicit-function result dU/dX|_S = -(dS/dX)/(dS/dU) (Section 4).
\begin{equation}
    \boxed{\left(\frac{\partial P}{\partial T}\right)_{\!V}
    \left(\frac{\partial T}{\partial V}\right)_{\!P}
    \left(\frac{\partial V}{\partial P}\right)_{\!T}=-1,
    \qquad
    \left(\frac{\partial P}{\partial T}\right)_{\!V}=\frac{\alpha_P}{\kappa_T}.}
\end{equation}

\paragraph{Relations between the response functions.}
The response functions are not independent. Two exact relations connect
the heat capacities and the compressibilities.
% TODO: derive C_P - C_V from S(T, V(T, P)) and the second Maxwell relation;
%       derive C_P / C_V = kappa_T / kappa_S from the cyclic relation.
\begin{equation}
    \boxed{C_P-C_V=\frac{TV\alpha_P^2}{\kappa_T},\qquad
    \frac{C_P}{C_V}=\frac{\kappa_T}{\kappa_S}.}
\end{equation}

\paragraph{Compressibility from particle-number fluctuations.}
The fluctuation formula of the grand canonical ensemble gives the
compressibility a statistical meaning, as the energy fluctuation does for
$C_V$. \textcolor{red}{Here, $n\equiv N/V$ denotes the number density.}
% FROM: Var(N) = k_B T (dN/dmu)_{T,V} (Section 5) and the Gibbs--Duhem
%       relation at fixed T, N dmu = V dP (Section 5).
\begin{equation}
    \boxed{\frac{\mathrm{Var}(N)}{N}=n\,k_BT\,\kappa_T.}
\end{equation}

% ============================================================
\subsection{Thermodynamic Stability}
% ============================================================

\paragraph{Positive heat capacities and compressibilities.}
Stability requires every response function that pairs a variable with the
conjugate of that variable to be positive.
% FROM: d^2S/dU^2 = -1/(C_V T^2) < 0 (Section 3); F convex in V (Section 4);
%       Var(N) >= 0 and the fluctuation formula above.
% TODO: show kappa_T > 0 from d^2F/dV^2 = 1/(V kappa_T) and from Var(N);
%       combine with the relations above to order C_P >= C_V and
%       kappa_T >= kappa_S.
\begin{equation}
    \boxed{C_P\geq C_V>0,\qquad\kappa_T\geq\kappa_S>0.}
\end{equation}

\paragraph{The principle of Le Chatelier.}
A stable system responds to an imposed change in a direction that reduces
the effect of the change. The positive response functions state this
principle quantitatively.
% TODO: one example each for heating at fixed V and compression at fixed T.

% ============================================================
\subsection{Thermodynamic Processes}
% ============================================================

\paragraph{Quasi-static processes and the path dependence of heat and work.}
A \textcolor{blue}{quasi-static process} passes through equilibrium states
only, so the state variables are defined along the whole path. Heat and work
then follow from integrals along the path.
% FROM: delta Q and delta W are path-dependent; delta Q <= T dS (Section 3).

\paragraph{Isothermal, isobaric, isochoric and adiabatic processes.}
Each elementary process fixes one variable. The heat and the work of each
process follow from the response functions.
% TODO: table of the four processes with delta Q, delta W and the response
%       function involved; adiabat versus isotherm in the P--V plane.
\begin{equation}
    \boxed{\left(\frac{\partial P}{\partial V}\right)_{\!S}
    =\frac{C_P}{C_V}\left(\frac{\partial P}{\partial V}\right)_{\!T}.}
\end{equation}

\paragraph{Free expansion.}
In a \textcolor{blue}{free expansion}, a gas expands into vacuum inside an
insulated container, so neither heat nor work is exchanged and $U$ stays
constant. The temperature change measures the deviation from ideal
behavior.
% TODO: derive with the second Maxwell relation; irreversible, so Delta S > 0.
\begin{equation}
    \boxed{\left(\frac{\partial T}{\partial V}\right)_{\!U}
    =-\frac{1}{C_V}\left[T\left(\frac{\partial P}{\partial T}\right)_{\!V}-P\right].}
\end{equation}

\paragraph{Throttling.}
In a \textcolor{blue}{throttling process}, a gas flows steadily through a
porous plug from high to low pressure. The enthalpy stays constant, and the
\textcolor{blue}{Joule--Thomson coefficient} gives the temperature change.
% TODO: show H = const from the work balance on both sides of the plug;
%       derive with the fourth Maxwell relation.
\begin{equation}
    \boxed{\mu_{\text{JT}}\equiv\left(\frac{\partial T}{\partial P}\right)_{\!H}
    =\frac{V}{C_P}\left(T\alpha_P-1\right).}
\end{equation}

\paragraph{Heat engines and the Carnot bound.}
A cyclic process returns to the initial state, so the net work equals the
net heat. The second law bounds the fraction of the absorbed heat that a
cycle converts into work.
% OPTIONAL: Carnot cycle; efficiency bound from the Clausius inequality.
\begin{equation}
    \boxed{\eta\leq1-\frac{T_{\text{low}}}{T_{\text{high}}}.}
\end{equation}

% ============================================================
\subsection{The Ideal Gas}
% ============================================================

\paragraph{Classical limit of the partition function.}
The sum over microstates becomes an integral over phase space for classical
particles. Two assumptions enter. Each microstate occupies a phase-space
volume $h^{3N}$, and the factor $1/N!$ removes the permutations of identical
particles.
% FROM: 1/N! and extensivity (Section 5, Gibbs paradox).
% TODO: state the h^{3N} cell as an assumption, justified later by quantum
%       statistics.
\begin{equation}
    \boxed{Z=\frac{1}{N!\,h^{3N}}\int d^{3N}q\,d^{3N}p\;e^{-\beta\mathcal{H}(q,p)}.}
\end{equation}

\paragraph{Partition function of non-interacting particles.}
The \textcolor{blue}{ideal gas} consists of point particles without
interactions, so the Hamiltonian contains the kinetic energy only. The
momentum integrals are Gaussian, and the position integrals give the
volume. \textcolor{red}{Here, $\lambda\equiv h/\sqrt{2\pi mk_BT}$ denotes the
thermal wavelength.}
\begin{equation}
    \boxed{Z=\frac{V^N}{N!\,\lambda^{3N}}.}
\end{equation}

\paragraph{Equation of state, energy and heat capacity.}
The free energy $F=-k_BT\ln Z$ gives the pressure, and
$U=-\partial\ln Z/\partial\beta$ gives the energy.
% FROM: F = -k_B T ln Z, P = -dF/dV (Section 4).
\begin{equation}
    \boxed{PV=Nk_BT,\qquad U=\tfrac{3}{2}Nk_BT,\qquad C_V=\tfrac{3}{2}Nk_B.}
\end{equation}

\paragraph{Entropy and chemical potential.}
The entropy of the ideal gas is the \textcolor{blue}{Sackur--Tetrode
equation}. The entropy is extensive only with the factor $1/N!$.
\begin{equation}
    \boxed{S=Nk_B\left[\ln\frac{V}{N\lambda^3}+\frac{5}{2}\right],\qquad
    \mu=k_BT\ln\!\left(n\lambda^3\right).}
\end{equation}

\paragraph{Response functions and processes of the ideal gas.}
The general results of this section take simple forms for the ideal gas.
Free expansion and throttling leave the temperature of an ideal gas
unchanged, because the particles do not interact.
% TODO: kappa_S, adiabat P V^gamma = const with gamma = 5/3; Var(N)/N = 1.
\begin{equation}
    \boxed{\alpha_P=\frac{1}{T},\qquad\kappa_T=\frac{1}{P},\qquad
    C_P-C_V=Nk_B,\qquad\mu_{\text{JT}}=0.}
\end{equation}

\paragraph{Validity of the ideal gas.}
The ideal gas requires two conditions. The gas must be dilute, so that the
interactions are negligible, and the thermal wavelength must be small
compared with the particle spacing, $n\lambda^3\ll1$, so that the classical
limit holds.

% ============================================================
\subsection{Interacting Gases}
% ============================================================

\paragraph{The configurational integral.}
Pair interactions $\phi(r)$ leave the momentum integrals unchanged. The
partition function then splits into the ideal-gas part and a
\textcolor{blue}{configurational integral} $Q_N$.
\begin{equation}
    \boxed{Z=\frac{Q_N}{N!\,\lambda^{3N}},\qquad
    Q_N=\int d^{3N}r\;\exp\!\Big[-\beta\sum_{i<j}\phi(r_{ij})\Big].}
\end{equation}

\paragraph{The virial expansion.}
At low density, the equation of state is a power series in $1/v$. The
\textcolor{blue}{second virial coefficient} $B_2(T)$ follows exactly from
the pair potential through the \textcolor{blue}{Mayer function}
$f(r)=e^{-\beta\phi(r)}-1$. This expansion is the rigorous link between the
interactions and the deviation from ideal behavior.
% TODO: derive B_2 from Q_N to first order in the density.
\begin{equation}
    \boxed{\frac{Pv}{k_BT}=1+\frac{B_2(T)}{v}+\cdots,\qquad
    B_2(T)=-\frac{1}{2}\int d^3r\,f(r).}
\end{equation}

\paragraph{The van der Waals equation.}
Two assumptions turn $Q_N$ into a closed form. A hard core excludes the
volume $b$ per particle, and the attraction acts on each particle as a
uniform mean field proportional to the density, with strength $a$.
% TODO: Q_N ~ (V - Nb)^N exp(beta a N^2 / V); compare B_2 = b - a/(k_B T)
%       with the virial result.
\begin{equation}
    \boxed{P=\frac{k_BT}{v-b}-\frac{a}{v^2}.}
\end{equation}

\paragraph{Critical point and corresponding states.}
Below a critical temperature, each van der Waals isotherm has a loop, and
the middle part of the loop has $\kappa_T<0$, which violates stability. The loop shrinks to a horizontal
inflection at the critical point, and the reduced variables
$P_r=P/P_c$, $v_r=v/v_c$ and $T_r=T/T_c$ remove the material constants.
\begin{equation}
    \boxed{v_c=3b,\quad k_BT_c=\frac{8a}{27b},\quad P_c=\frac{a}{27b^2},
    \qquad
    \left(P_r+\frac{3}{v_r^2}\right)\left(3v_r-1\right)=8T_r.}
\end{equation}
% TODO: law of corresponding states; the loop is resolved by phase
%       coexistence in the phase-transition sections (Maxwell construction).

\paragraph{The Dieterici equation.}
The \textcolor{blue}{Dieterici equation} replaces the mean-field attraction
by a Boltzmann factor and has the same second virial coefficient as the
van der Waals equation. The critical compressibility factor
$Z_c=P_cv_c/(k_BT_c)$ tests both models, because $Z_c$ is the same for all
substances in either model. Table~\ref{tab:Zc} compares the models with
experiment.
\begin{equation}
    \boxed{P=\frac{k_BT}{v-b}\exp\!\left(-\frac{a}{k_BTv}\right).}
\end{equation}
\begin{table}[h]
\centering
\caption{Critical compressibility factor $Z_c$ of two model equations and
four fluids.}
\label{tab:Zc}
\begin{tabular}{lc}
    Equation or fluid & $Z_c$ \\ \hline
    van der Waals & $3/8=0.375$ \\
    Dieterici & $2/e^2\approx0.271$ \\
    Argon & 0.291 \\
    Nitrogen & 0.289 \\
    Carbon dioxide & 0.274 \\
    Water & 0.229 \\
\end{tabular}
\end{table}

\paragraph{Joule--Thomson inversion.}
The attraction cools a throttled gas, and the hard core heats the gas. At
low density, the van der Waals gas cools on throttling below the
\textcolor{blue}{inversion temperature}.
% TODO: derive mu_JT ~ (2a/(k_B T) - b)/c_P from B_2.
\begin{equation}
    \boxed{k_BT_{\text{inv}}=\frac{2a}{b}=\frac{27}{4}\,k_BT_c.}
\end{equation}

% under \part{Classical Thermodynamics} and move that part in front of
% \part{Phase Transitions}.
%
% Conventions of this skeleton:
%   - Each \paragraph has one or two sentences that state the goal of the
%     paragraph, followed by the boxed target result where one exists.
%   - "% TODO" comments list the derivation steps still to be written.
%   - "% FROM" comments name the earlier result that a step reuses, so the
%     chapter does not repeat Sections 1-5.
% ============================================================

\section{Classical Thermodynamics of Simple Fluids}

\paragraph{Scope of this section.}
Sections~3--5 derive the laws of thermodynamics, the potentials and the
fluctuation formulas from entropy maximization. This section turns these
results into measurable quantities and applies them to model fluids. The
response functions and the stability conditions of this section are the
tools of the phase-transition sections, and the ideal gas and the van der
Waals gas are the first explicit models of the notes.

% ============================================================
\subsection{Response Functions}
% ============================================================

\paragraph{Measurable second derivatives.}
Experiments measure how one variable responds to a change of another,
not the potentials themselves. Each \textcolor{blue}{response function} is a
second derivative of a potential.
% FROM: C_V and Var(U) = k_B T^2 C_V (Section 3); delta Q = dH at fixed P
%       (Section 4); G(T, P, N) and its derivatives (Section 5).
% TODO: define C_P = T (dS/dT)_P = (dH/dT)_P, alpha_P, kappa_T,
%       kappa_S = -(1/V)(dV/dP)_S; write each as a second derivative of G.
\begin{equation}
    \boxed{C_P=-T\frac{\partial^2G}{\partial T^2},\qquad
    \alpha_P=\frac{1}{V}\frac{\partial^2G}{\partial T\,\partial P},\qquad
    \kappa_T=-\frac{1}{V}\frac{\partial^2G}{\partial P^2}.}
\end{equation}

\paragraph{Maxwell relations.}
The mixed second derivatives of a potential do not depend on the order of
differentiation. Each potential therefore gives one
\textcolor{blue}{Maxwell relation} between two response-type derivatives.
% TODO: derive from dU, dF, dH, dG (Sections 3-5); note which relation turns
%       an unmeasurable entropy derivative into a measurable one.
\begin{equation}
    \boxed{
    \begin{aligned}
        \left(\frac{\partial T}{\partial V}\right)_{\!S,N}&=-\left(\frac{\partial P}{\partial S}\right)_{\!V,N}, &
        \left(\frac{\partial S}{\partial V}\right)_{\!T,N}&=\left(\frac{\partial P}{\partial T}\right)_{\!V,N},\\
        \left(\frac{\partial T}{\partial P}\right)_{\!S,N}&=\left(\frac{\partial V}{\partial S}\right)_{\!P,N}, &
        \left(\frac{\partial S}{\partial P}\right)_{\!T,N}&=-\left(\frac{\partial V}{\partial T}\right)_{\!P,N}.
    \end{aligned}}
\end{equation}

\paragraph{The cyclic relation.}
Three variables linked by one equation of state satisfy the cyclic
relation. The slope of an isochore in the $P$--$T$ plane then follows from
$\alpha_P$ and $\kappa_T$.
% FROM: the implicit-function result dU/dX|_S = -(dS/dX)/(dS/dU) (Section 4).
\begin{equation}
    \boxed{\left(\frac{\partial P}{\partial T}\right)_{\!V}
    \left(\frac{\partial T}{\partial V}\right)_{\!P}
    \left(\frac{\partial V}{\partial P}\right)_{\!T}=-1,
    \qquad
    \left(\frac{\partial P}{\partial T}\right)_{\!V}=\frac{\alpha_P}{\kappa_T}.}
\end{equation}

\paragraph{Relations between the response functions.}
The response functions are not independent. Two exact relations connect
the heat capacities and the compressibilities.
% TODO: derive C_P - C_V from S(T, V(T, P)) and the second Maxwell relation;
%       derive C_P / C_V = kappa_T / kappa_S from the cyclic relation.
\begin{equation}
    \boxed{C_P-C_V=\frac{TV\alpha_P^2}{\kappa_T},\qquad
    \frac{C_P}{C_V}=\frac{\kappa_T}{\kappa_S}.}
\end{equation}

\paragraph{Compressibility from particle-number fluctuations.}
The fluctuation formula of the grand canonical ensemble gives the
compressibility a statistical meaning, as the energy fluctuation does for
$C_V$. \textcolor{red}{Here, $n\equiv N/V$ denotes the number density.}
% FROM: Var(N) = k_B T (dN/dmu)_{T,V} (Section 5) and the Gibbs--Duhem
%       relation at fixed T, N dmu = V dP (Section 5).
\begin{equation}
    \boxed{\frac{\mathrm{Var}(N)}{N}=n\,k_BT\,\kappa_T.}
\end{equation}

% ============================================================
\subsection{Thermodynamic Stability}
% ============================================================

\paragraph{Positive heat capacities and compressibilities.}
Stability requires every response function that pairs a variable with the
conjugate of that variable to be positive.
% FROM: d^2S/dU^2 = -1/(C_V T^2) < 0 (Section 3); F convex in V (Section 4);
%       Var(N) >= 0 and the fluctuation formula above.
% TODO: show kappa_T > 0 from d^2F/dV^2 = 1/(V kappa_T) and from Var(N);
%       combine with the relations above to order C_P >= C_V and
%       kappa_T >= kappa_S.
\begin{equation}
    \boxed{C_P\geq C_V>0,\qquad\kappa_T\geq\kappa_S>0.}
\end{equation}

\paragraph{The principle of Le Chatelier.}
A stable system responds to an imposed change in a direction that reduces
the effect of the change. The positive response functions state this
principle quantitatively.
% TODO: one example each for heating at fixed V and compression at fixed T.

% ============================================================
\subsection{Thermodynamic Processes}
% ============================================================

\paragraph{Quasi-static processes and the path dependence of heat and work.}
A \textcolor{blue}{quasi-static process} passes through equilibrium states
only, so the state variables are defined along the whole path. Heat and work
then follow from integrals along the path.
% FROM: delta Q and delta W are path-dependent; delta Q <= T dS (Section 3).

\paragraph{Isothermal, isobaric, isochoric and adiabatic processes.}
Each elementary process fixes one variable. The heat and the work of each
process follow from the response functions.
% TODO: table of the four processes with delta Q, delta W and the response
%       function involved; adiabat versus isotherm in the P--V plane.
\begin{equation}
    \boxed{\left(\frac{\partial P}{\partial V}\right)_{\!S}
    =\frac{C_P}{C_V}\left(\frac{\partial P}{\partial V}\right)_{\!T}.}
\end{equation}

\paragraph{Free expansion.}
In a \textcolor{blue}{free expansion}, a gas expands into vacuum inside an
insulated container, so neither heat nor work is exchanged and $U$ stays
constant. The temperature change measures the deviation from ideal
behavior.
% TODO: derive with the second Maxwell relation; irreversible, so Delta S > 0.
\begin{equation}
    \boxed{\left(\frac{\partial T}{\partial V}\right)_{\!U}
    =-\frac{1}{C_V}\left[T\left(\frac{\partial P}{\partial T}\right)_{\!V}-P\right].}
\end{equation}

\paragraph{Throttling.}
In a \textcolor{blue}{throttling process}, a gas flows steadily through a
porous plug from high to low pressure. The enthalpy stays constant, and the
\textcolor{blue}{Joule--Thomson coefficient} gives the temperature change.
% TODO: show H = const from the work balance on both sides of the plug;
%       derive with the fourth Maxwell relation.
\begin{equation}
    \boxed{\mu_{\text{JT}}\equiv\left(\frac{\partial T}{\partial P}\right)_{\!H}
    =\frac{V}{C_P}\left(T\alpha_P-1\right).}
\end{equation}

\paragraph{Heat engines and the Carnot bound.}
A cyclic process returns to the initial state, so the net work equals the
net heat. The second law bounds the fraction of the absorbed heat that a
cycle converts into work.
% OPTIONAL: Carnot cycle; efficiency bound from the Clausius inequality.
\begin{equation}
    \boxed{\eta\leq1-\frac{T_{\text{low}}}{T_{\text{high}}}.}
\end{equation}

% ============================================================
\subsection{The Ideal Gas}
% ============================================================

\paragraph{Classical limit of the partition function.}
The sum over microstates becomes an integral over phase space for classical
particles. Two assumptions enter. Each microstate occupies a phase-space
volume $h^{3N}$, and the factor $1/N!$ removes the permutations of identical
particles.
% FROM: 1/N! and extensivity (Section 5, Gibbs paradox).
% TODO: state the h^{3N} cell as an assumption, justified later by quantum
%       statistics.
\begin{equation}
    \boxed{Z=\frac{1}{N!\,h^{3N}}\int d^{3N}q\,d^{3N}p\;e^{-\beta\mathcal{H}(q,p)}.}
\end{equation}

\paragraph{Partition function of non-interacting particles.}
The \textcolor{blue}{ideal gas} consists of point particles without
interactions, so the Hamiltonian contains the kinetic energy only. The
momentum integrals are Gaussian, and the position integrals give the
volume. \textcolor{red}{Here, $\lambda\equiv h/\sqrt{2\pi mk_BT}$ denotes the
thermal wavelength.}
\begin{equation}
    \boxed{Z=\frac{V^N}{N!\,\lambda^{3N}}.}
\end{equation}

\paragraph{Equation of state, energy and heat capacity.}
The free energy $F=-k_BT\ln Z$ gives the pressure, and
$U=-\partial\ln Z/\partial\beta$ gives the energy.
% FROM: F = -k_B T ln Z, P = -dF/dV (Section 4).
\begin{equation}
    \boxed{PV=Nk_BT,\qquad U=\tfrac{3}{2}Nk_BT,\qquad C_V=\tfrac{3}{2}Nk_B.}
\end{equation}

\paragraph{Entropy and chemical potential.}
The entropy of the ideal gas is the \textcolor{blue}{Sackur--Tetrode
equation}. The entropy is extensive only with the factor $1/N!$.
\begin{equation}
    \boxed{S=Nk_B\left[\ln\frac{V}{N\lambda^3}+\frac{5}{2}\right],\qquad
    \mu=k_BT\ln\!\left(n\lambda^3\right).}
\end{equation}

\paragraph{Response functions and processes of the ideal gas.}
The general results of this section take simple forms for the ideal gas.
Free expansion and throttling leave the temperature of an ideal gas
unchanged, because the particles do not interact.
% TODO: kappa_S, adiabat P V^gamma = const with gamma = 5/3; Var(N)/N = 1.
\begin{equation}
    \boxed{\alpha_P=\frac{1}{T},\qquad\kappa_T=\frac{1}{P},\qquad
    C_P-C_V=Nk_B,\qquad\mu_{\text{JT}}=0.}
\end{equation}

\paragraph{Validity of the ideal gas.}
The ideal gas requires two conditions. The gas must be dilute, so that the
interactions are negligible, and the thermal wavelength must be small
compared with the particle spacing, $n\lambda^3\ll1$, so that the classical
limit holds.

% ============================================================
\subsection{Interacting Gases}
% ============================================================

\paragraph{The configurational integral.}
Pair interactions $\phi(r)$ leave the momentum integrals unchanged. The
partition function then splits into the ideal-gas part and a
\textcolor{blue}{configurational integral} $Q_N$.
\begin{equation}
    \boxed{Z=\frac{Q_N}{N!\,\lambda^{3N}},\qquad
    Q_N=\int d^{3N}r\;\exp\!\Big[-\beta\sum_{i<j}\phi(r_{ij})\Big].}
\end{equation}

\paragraph{The virial expansion.}
At low density, the equation of state is a power series in $1/v$. The
\textcolor{blue}{second virial coefficient} $B_2(T)$ follows exactly from
the pair potential through the \textcolor{blue}{Mayer function}
$f(r)=e^{-\beta\phi(r)}-1$. This expansion is the rigorous link between the
interactions and the deviation from ideal behavior.
% TODO: derive B_2 from Q_N to first order in the density.
\begin{equation}
    \boxed{\frac{Pv}{k_BT}=1+\frac{B_2(T)}{v}+\cdots,\qquad
    B_2(T)=-\frac{1}{2}\int d^3r\,f(r).}
\end{equation}

\paragraph{The van der Waals equation.}
Two assumptions turn $Q_N$ into a closed form. A hard core excludes the
volume $b$ per particle, and the attraction acts on each particle as a
uniform mean field proportional to the density, with strength $a$.
% TODO: Q_N ~ (V - Nb)^N exp(beta a N^2 / V); compare B_2 = b - a/(k_B T)
%       with the virial result.
\begin{equation}
    \boxed{P=\frac{k_BT}{v-b}-\frac{a}{v^2}.}
\end{equation}

\paragraph{Critical point and corresponding states.}
Below a critical temperature, each van der Waals isotherm has a loop, and
the middle part of the loop has $\kappa_T<0$, which violates stability. The loop shrinks to a horizontal
inflection at the critical point, and the reduced variables
$P_r=P/P_c$, $v_r=v/v_c$ and $T_r=T/T_c$ remove the material constants.
\begin{equation}
    \boxed{v_c=3b,\quad k_BT_c=\frac{8a}{27b},\quad P_c=\frac{a}{27b^2},
    \qquad
    \left(P_r+\frac{3}{v_r^2}\right)\left(3v_r-1\right)=8T_r.}
\end{equation}
% TODO: law of corresponding states; the loop is resolved by phase
%       coexistence in the phase-transition sections (Maxwell construction).

\paragraph{The Dieterici equation.}
The \textcolor{blue}{Dieterici equation} replaces the mean-field attraction
by a Boltzmann factor and has the same second virial coefficient as the
van der Waals equation. The critical compressibility factor
$Z_c=P_cv_c/(k_BT_c)$ tests both models, because $Z_c$ is the same for all
substances in either model. Table~\ref{tab:Zc} compares the models with
experiment.
\begin{equation}
    \boxed{P=\frac{k_BT}{v-b}\exp\!\left(-\frac{a}{k_BTv}\right).}
\end{equation}
\begin{table}[h]
\centering
\caption{Critical compressibility factor $Z_c$ of two model equations and
four fluids.}
\label{tab:Zc}
\begin{tabular}{lc}
    Equation or fluid & $Z_c$ \\ \hline
    van der Waals & $3/8=0.375$ \\
    Dieterici & $2/e^2\approx0.271$ \\
    Argon & 0.291 \\
    Nitrogen & 0.289 \\
    Carbon dioxide & 0.274 \\
    Water & 0.229 \\
\end{tabular}
\end{table}

\paragraph{Joule--Thomson inversion.}
The attraction cools a throttled gas, and the hard core heats the gas. At
low density, the van der Waals gas cools on throttling below the
\textcolor{blue}{inversion temperature}.
% TODO: derive mu_JT ~ (2a/(k_B T) - b)/c_P from B_2.
\begin{equation}
    \boxed{k_BT_{\text{inv}}=\frac{2a}{b}=\frac{27}{4}\,k_BT_c.}
\end{equation}

% under \part{Classical Thermodynamics} and move that part in front of
% \part{Phase Transitions}.
%
% Conventions of this skeleton:
%   - Each \paragraph has one or two sentences that state the goal of the
%     paragraph, followed by the boxed target result where one exists.
%   - "% TODO" comments list the derivation steps still to be written.
%   - "% FROM" comments name the earlier result that a step reuses, so the
%     chapter does not repeat Sections 1-5.
% ============================================================

\section{Classical Thermodynamics of Simple Fluids}

\paragraph{Scope of this section.}
Sections~3--5 derive the laws of thermodynamics, the potentials and the
fluctuation formulas from entropy maximization. This section turns these
results into measurable quantities and applies them to model fluids. The
response functions and the stability conditions of this section are the
tools of the phase-transition sections, and the ideal gas and the van der
Waals gas are the first explicit models of the notes.

% ============================================================
\subsection{Response Functions}
% ============================================================

\paragraph{Measurable second derivatives.}
Experiments measure how one variable responds to a change of another,
not the potentials themselves. Each \textcolor{blue}{response function} is a
second derivative of a potential.
% FROM: C_V and Var(U) = k_B T^2 C_V (Section 3); delta Q = dH at fixed P
%       (Section 4); G(T, P, N) and its derivatives (Section 5).
% TODO: define C_P = T (dS/dT)_P = (dH/dT)_P, alpha_P, kappa_T,
%       kappa_S = -(1/V)(dV/dP)_S; write each as a second derivative of G.
\begin{equation}
    \boxed{C_P=-T\frac{\partial^2G}{\partial T^2},\qquad
    \alpha_P=\frac{1}{V}\frac{\partial^2G}{\partial T\,\partial P},\qquad
    \kappa_T=-\frac{1}{V}\frac{\partial^2G}{\partial P^2}.}
\end{equation}

\paragraph{Maxwell relations.}
The mixed second derivatives of a potential do not depend on the order of
differentiation. Each potential therefore gives one
\textcolor{blue}{Maxwell relation} between two response-type derivatives.
% TODO: derive from dU, dF, dH, dG (Sections 3-5); note which relation turns
%       an unmeasurable entropy derivative into a measurable one.
\begin{equation}
    \boxed{
    \begin{aligned}
        \left(\frac{\partial T}{\partial V}\right)_{\!S,N}&=-\left(\frac{\partial P}{\partial S}\right)_{\!V,N}, &
        \left(\frac{\partial S}{\partial V}\right)_{\!T,N}&=\left(\frac{\partial P}{\partial T}\right)_{\!V,N},\\
        \left(\frac{\partial T}{\partial P}\right)_{\!S,N}&=\left(\frac{\partial V}{\partial S}\right)_{\!P,N}, &
        \left(\frac{\partial S}{\partial P}\right)_{\!T,N}&=-\left(\frac{\partial V}{\partial T}\right)_{\!P,N}.
    \end{aligned}}
\end{equation}

\paragraph{The cyclic relation.}
Three variables linked by one equation of state satisfy the cyclic
relation. The slope of an isochore in the $P$--$T$ plane then follows from
$\alpha_P$ and $\kappa_T$.
% FROM: the implicit-function result dU/dX|_S = -(dS/dX)/(dS/dU) (Section 4).
\begin{equation}
    \boxed{\left(\frac{\partial P}{\partial T}\right)_{\!V}
    \left(\frac{\partial T}{\partial V}\right)_{\!P}
    \left(\frac{\partial V}{\partial P}\right)_{\!T}=-1,
    \qquad
    \left(\frac{\partial P}{\partial T}\right)_{\!V}=\frac{\alpha_P}{\kappa_T}.}
\end{equation}

\paragraph{Relations between the response functions.}
The response functions are not independent. Two exact relations connect
the heat capacities and the compressibilities.
% TODO: derive C_P - C_V from S(T, V(T, P)) and the second Maxwell relation;
%       derive C_P / C_V = kappa_T / kappa_S from the cyclic relation.
\begin{equation}
    \boxed{C_P-C_V=\frac{TV\alpha_P^2}{\kappa_T},\qquad
    \frac{C_P}{C_V}=\frac{\kappa_T}{\kappa_S}.}
\end{equation}

\paragraph{Compressibility from particle-number fluctuations.}
The fluctuation formula of the grand canonical ensemble gives the
compressibility a statistical meaning, as the energy fluctuation does for
$C_V$. \textcolor{red}{Here, $n\equiv N/V$ denotes the number density.}
% FROM: Var(N) = k_B T (dN/dmu)_{T,V} (Section 5) and the Gibbs--Duhem
%       relation at fixed T, N dmu = V dP (Section 5).
\begin{equation}
    \boxed{\frac{\mathrm{Var}(N)}{N}=n\,k_BT\,\kappa_T.}
\end{equation}

% ============================================================
\subsection{Thermodynamic Stability}
% ============================================================

\paragraph{Positive heat capacities and compressibilities.}
Stability requires every response function that pairs a variable with the
conjugate of that variable to be positive.
% FROM: d^2S/dU^2 = -1/(C_V T^2) < 0 (Section 3); F convex in V (Section 4);
%       Var(N) >= 0 and the fluctuation formula above.
% TODO: show kappa_T > 0 from d^2F/dV^2 = 1/(V kappa_T) and from Var(N);
%       combine with the relations above to order C_P >= C_V and
%       kappa_T >= kappa_S.
\begin{equation}
    \boxed{C_P\geq C_V>0,\qquad\kappa_T\geq\kappa_S>0.}
\end{equation}

\paragraph{The principle of Le Chatelier.}
A stable system responds to an imposed change in a direction that reduces
the effect of the change. The positive response functions state this
principle quantitatively.
% TODO: one example each for heating at fixed V and compression at fixed T.

% ============================================================
\subsection{Thermodynamic Processes}
% ============================================================

\paragraph{Quasi-static processes and the path dependence of heat and work.}
A \textcolor{blue}{quasi-static process} passes through equilibrium states
only, so the state variables are defined along the whole path. Heat and work
then follow from integrals along the path.
% FROM: delta Q and delta W are path-dependent; delta Q <= T dS (Section 3).

\paragraph{Isothermal, isobaric, isochoric and adiabatic processes.}
Each elementary process fixes one variable. The heat and the work of each
process follow from the response functions.
% TODO: table of the four processes with delta Q, delta W and the response
%       function involved; adiabat versus isotherm in the P--V plane.
\begin{equation}
    \boxed{\left(\frac{\partial P}{\partial V}\right)_{\!S}
    =\frac{C_P}{C_V}\left(\frac{\partial P}{\partial V}\right)_{\!T}.}
\end{equation}

\paragraph{Free expansion.}
In a \textcolor{blue}{free expansion}, a gas expands into vacuum inside an
insulated container, so neither heat nor work is exchanged and $U$ stays
constant. The temperature change measures the deviation from ideal
behavior.
% TODO: derive with the second Maxwell relation; irreversible, so Delta S > 0.
\begin{equation}
    \boxed{\left(\frac{\partial T}{\partial V}\right)_{\!U}
    =-\frac{1}{C_V}\left[T\left(\frac{\partial P}{\partial T}\right)_{\!V}-P\right].}
\end{equation}

\paragraph{Throttling.}
In a \textcolor{blue}{throttling process}, a gas flows steadily through a
porous plug from high to low pressure. The enthalpy stays constant, and the
\textcolor{blue}{Joule--Thomson coefficient} gives the temperature change.
% TODO: show H = const from the work balance on both sides of the plug;
%       derive with the fourth Maxwell relation.
\begin{equation}
    \boxed{\mu_{\text{JT}}\equiv\left(\frac{\partial T}{\partial P}\right)_{\!H}
    =\frac{V}{C_P}\left(T\alpha_P-1\right).}
\end{equation}

\paragraph{Heat engines and the Carnot bound.}
A cyclic process returns to the initial state, so the net work equals the
net heat. The second law bounds the fraction of the absorbed heat that a
cycle converts into work.
% OPTIONAL: Carnot cycle; efficiency bound from the Clausius inequality.
\begin{equation}
    \boxed{\eta\leq1-\frac{T_{\text{low}}}{T_{\text{high}}}.}
\end{equation}

% ============================================================
\subsection{The Ideal Gas}
% ============================================================

\paragraph{Classical limit of the partition function.}
The sum over microstates becomes an integral over phase space for classical
particles. Two assumptions enter. Each microstate occupies a phase-space
volume $h^{3N}$, and the factor $1/N!$ removes the permutations of identical
particles.
% FROM: 1/N! and extensivity (Section 5, Gibbs paradox).
% TODO: state the h^{3N} cell as an assumption, justified later by quantum
%       statistics.
\begin{equation}
    \boxed{Z=\frac{1}{N!\,h^{3N}}\int d^{3N}q\,d^{3N}p\;e^{-\beta\mathcal{H}(q,p)}.}
\end{equation}

\paragraph{Partition function of non-interacting particles.}
The \textcolor{blue}{ideal gas} consists of point particles without
interactions, so the Hamiltonian contains the kinetic energy only. The
momentum integrals are Gaussian, and the position integrals give the
volume. \textcolor{red}{Here, $\lambda\equiv h/\sqrt{2\pi mk_BT}$ denotes the
thermal wavelength.}
\begin{equation}
    \boxed{Z=\frac{V^N}{N!\,\lambda^{3N}}.}
\end{equation}

\paragraph{Equation of state, energy and heat capacity.}
The free energy $F=-k_BT\ln Z$ gives the pressure, and
$U=-\partial\ln Z/\partial\beta$ gives the energy.
% FROM: F = -k_B T ln Z, P = -dF/dV (Section 4).
\begin{equation}
    \boxed{PV=Nk_BT,\qquad U=\tfrac{3}{2}Nk_BT,\qquad C_V=\tfrac{3}{2}Nk_B.}
\end{equation}

\paragraph{Entropy and chemical potential.}
The entropy of the ideal gas is the \textcolor{blue}{Sackur--Tetrode
equation}. The entropy is extensive only with the factor $1/N!$.
\begin{equation}
    \boxed{S=Nk_B\left[\ln\frac{V}{N\lambda^3}+\frac{5}{2}\right],\qquad
    \mu=k_BT\ln\!\left(n\lambda^3\right).}
\end{equation}

\paragraph{Response functions and processes of the ideal gas.}
The general results of this section take simple forms for the ideal gas.
Free expansion and throttling leave the temperature of an ideal gas
unchanged, because the particles do not interact.
% TODO: kappa_S, adiabat P V^gamma = const with gamma = 5/3; Var(N)/N = 1.
\begin{equation}
    \boxed{\alpha_P=\frac{1}{T},\qquad\kappa_T=\frac{1}{P},\qquad
    C_P-C_V=Nk_B,\qquad\mu_{\text{JT}}=0.}
\end{equation}

\paragraph{Validity of the ideal gas.}
The ideal gas requires two conditions. The gas must be dilute, so that the
interactions are negligible, and the thermal wavelength must be small
compared with the particle spacing, $n\lambda^3\ll1$, so that the classical
limit holds.

% ============================================================
\subsection{Interacting Gases}
% ============================================================

\paragraph{The configurational integral.}
Pair interactions $\phi(r)$ leave the momentum integrals unchanged. The
partition function then splits into the ideal-gas part and a
\textcolor{blue}{configurational integral} $Q_N$.
\begin{equation}
    \boxed{Z=\frac{Q_N}{N!\,\lambda^{3N}},\qquad
    Q_N=\int d^{3N}r\;\exp\!\Big[-\beta\sum_{i<j}\phi(r_{ij})\Big].}
\end{equation}

\paragraph{The virial expansion.}
At low density, the equation of state is a power series in $1/v$. The
\textcolor{blue}{second virial coefficient} $B_2(T)$ follows exactly from
the pair potential through the \textcolor{blue}{Mayer function}
$f(r)=e^{-\beta\phi(r)}-1$. This expansion is the rigorous link between the
interactions and the deviation from ideal behavior.
% TODO: derive B_2 from Q_N to first order in the density.
\begin{equation}
    \boxed{\frac{Pv}{k_BT}=1+\frac{B_2(T)}{v}+\cdots,\qquad
    B_2(T)=-\frac{1}{2}\int d^3r\,f(r).}
\end{equation}

\paragraph{The van der Waals equation.}
Two assumptions turn $Q_N$ into a closed form. A hard core excludes the
volume $b$ per particle, and the attraction acts on each particle as a
uniform mean field proportional to the density, with strength $a$.
% TODO: Q_N ~ (V - Nb)^N exp(beta a N^2 / V); compare B_2 = b - a/(k_B T)
%       with the virial result.
\begin{equation}
    \boxed{P=\frac{k_BT}{v-b}-\frac{a}{v^2}.}
\end{equation}

\paragraph{Critical point and corresponding states.}
Below a critical temperature, each van der Waals isotherm has a loop, and
the middle part of the loop has $\kappa_T<0$, which violates stability. The loop shrinks to a horizontal
inflection at the critical point, and the reduced variables
$P_r=P/P_c$, $v_r=v/v_c$ and $T_r=T/T_c$ remove the material constants.
\begin{equation}
    \boxed{v_c=3b,\quad k_BT_c=\frac{8a}{27b},\quad P_c=\frac{a}{27b^2},
    \qquad
    \left(P_r+\frac{3}{v_r^2}\right)\left(3v_r-1\right)=8T_r.}
\end{equation}
% TODO: law of corresponding states; the loop is resolved by phase
%       coexistence in the phase-transition sections (Maxwell construction).

\paragraph{The Dieterici equation.}
The \textcolor{blue}{Dieterici equation} replaces the mean-field attraction
by a Boltzmann factor and has the same second virial coefficient as the
van der Waals equation. The critical compressibility factor
$Z_c=P_cv_c/(k_BT_c)$ tests both models, because $Z_c$ is the same for all
substances in either model. Table~\ref{tab:Zc} compares the models with
experiment.
\begin{equation}
    \boxed{P=\frac{k_BT}{v-b}\exp\!\left(-\frac{a}{k_BTv}\right).}
\end{equation}
\begin{table}[h]
\centering
\caption{Critical compressibility factor $Z_c$ of two model equations and
four fluids.}
\label{tab:Zc}
\begin{tabular}{lc}
    Equation or fluid & $Z_c$ \\ \hline
    van der Waals & $3/8=0.375$ \\
    Dieterici & $2/e^2\approx0.271$ \\
    Argon & 0.291 \\
    Nitrogen & 0.289 \\
    Carbon dioxide & 0.274 \\
    Water & 0.229 \\
\end{tabular}
\end{table}

\paragraph{Joule--Thomson inversion.}
The attraction cools a throttled gas, and the hard core heats the gas. At
low density, the van der Waals gas cools on throttling below the
\textcolor{blue}{inversion temperature}.
% TODO: derive mu_JT ~ (2a/(k_B T) - b)/c_P from B_2.
\begin{equation}
    \boxed{k_BT_{\text{inv}}=\frac{2a}{b}=\frac{27}{4}\,k_BT_c.}
\end{equation}

% under \part{Classical Thermodynamics} and move that part in front of
% \part{Phase Transitions}.
%
% Conventions of this skeleton:
%   - Each \paragraph has one or two sentences that state the goal of the
%     paragraph, followed by the boxed target result where one exists.
%   - "% TODO" comments list the derivation steps still to be written.
%   - "% FROM" comments name the earlier result that a step reuses, so the
%     chapter does not repeat Sections 1-5.
% ============================================================

\section{Classical Thermodynamics of Simple Fluids}

\paragraph{Scope of this section.}
Sections~3--5 derive the laws of thermodynamics, the potentials and the
fluctuation formulas from entropy maximization. This section turns these
results into measurable quantities and applies them to model fluids. The
response functions and the stability conditions of this section are the
tools of the phase-transition sections, and the ideal gas and the van der
Waals gas are the first explicit models of the notes.

% ============================================================
\subsection{Response Functions}
% ============================================================

\paragraph{Measurable second derivatives.}
Experiments measure how one variable responds to a change of another,
not the potentials themselves. Each \textcolor{blue}{response function} is a
second derivative of a potential.
% FROM: C_V and Var(U) = k_B T^2 C_V (Section 3); delta Q = dH at fixed P
%       (Section 4); G(T, P, N) and its derivatives (Section 5).
% TODO: define C_P = T (dS/dT)_P = (dH/dT)_P, alpha_P, kappa_T,
%       kappa_S = -(1/V)(dV/dP)_S; write each as a second derivative of G.
\begin{equation}
    \boxed{C_P=-T\frac{\partial^2G}{\partial T^2},\qquad
    \alpha_P=\frac{1}{V}\frac{\partial^2G}{\partial T\,\partial P},\qquad
    \kappa_T=-\frac{1}{V}\frac{\partial^2G}{\partial P^2}.}
\end{equation}

\paragraph{Maxwell relations.}
The mixed second derivatives of a potential do not depend on the order of
differentiation. Each potential therefore gives one
\textcolor{blue}{Maxwell relation} between two response-type derivatives.
% TODO: derive from dU, dF, dH, dG (Sections 3-5); note which relation turns
%       an unmeasurable entropy derivative into a measurable one.
\begin{equation}
    \boxed{
    \begin{aligned}
        \left(\frac{\partial T}{\partial V}\right)_{\!S,N}&=-\left(\frac{\partial P}{\partial S}\right)_{\!V,N}, &
        \left(\frac{\partial S}{\partial V}\right)_{\!T,N}&=\left(\frac{\partial P}{\partial T}\right)_{\!V,N},\\
        \left(\frac{\partial T}{\partial P}\right)_{\!S,N}&=\left(\frac{\partial V}{\partial S}\right)_{\!P,N}, &
        \left(\frac{\partial S}{\partial P}\right)_{\!T,N}&=-\left(\frac{\partial V}{\partial T}\right)_{\!P,N}.
    \end{aligned}}
\end{equation}

\paragraph{The cyclic relation.}
Three variables linked by one equation of state satisfy the cyclic
relation. The slope of an isochore in the $P$--$T$ plane then follows from
$\alpha_P$ and $\kappa_T$.
% FROM: the implicit-function result dU/dX|_S = -(dS/dX)/(dS/dU) (Section 4).
\begin{equation}
    \boxed{\left(\frac{\partial P}{\partial T}\right)_{\!V}
    \left(\frac{\partial T}{\partial V}\right)_{\!P}
    \left(\frac{\partial V}{\partial P}\right)_{\!T}=-1,
    \qquad
    \left(\frac{\partial P}{\partial T}\right)_{\!V}=\frac{\alpha_P}{\kappa_T}.}
\end{equation}

\paragraph{Relations between the response functions.}
The response functions are not independent. Two exact relations connect
the heat capacities and the compressibilities.
% TODO: derive C_P - C_V from S(T, V(T, P)) and the second Maxwell relation;
%       derive C_P / C_V = kappa_T / kappa_S from the cyclic relation.
\begin{equation}
    \boxed{C_P-C_V=\frac{TV\alpha_P^2}{\kappa_T},\qquad
    \frac{C_P}{C_V}=\frac{\kappa_T}{\kappa_S}.}
\end{equation}

\paragraph{Compressibility from particle-number fluctuations.}
The fluctuation formula of the grand canonical ensemble gives the
compressibility a statistical meaning, as the energy fluctuation does for
$C_V$. \textcolor{red}{Here, $n\equiv N/V$ denotes the number density.}
% FROM: Var(N) = k_B T (dN/dmu)_{T,V} (Section 5) and the Gibbs--Duhem
%       relation at fixed T, N dmu = V dP (Section 5).
\begin{equation}
    \boxed{\frac{\mathrm{Var}(N)}{N}=n\,k_BT\,\kappa_T.}
\end{equation}

% ============================================================
\subsection{Thermodynamic Stability}
% ============================================================

\paragraph{Positive heat capacities and compressibilities.}
Stability requires every response function that pairs a variable with the
conjugate of that variable to be positive.
% FROM: d^2S/dU^2 = -1/(C_V T^2) < 0 (Section 3); F convex in V (Section 4);
%       Var(N) >= 0 and the fluctuation formula above.
% TODO: show kappa_T > 0 from d^2F/dV^2 = 1/(V kappa_T) and from Var(N);
%       combine with the relations above to order C_P >= C_V and
%       kappa_T >= kappa_S.
\begin{equation}
    \boxed{C_P\geq C_V>0,\qquad\kappa_T\geq\kappa_S>0.}
\end{equation}

\paragraph{The principle of Le Chatelier.}
A stable system responds to an imposed change in a direction that reduces
the effect of the change. The positive response functions state this
principle quantitatively.
% TODO: one example each for heating at fixed V and compression at fixed T.

% ============================================================
\subsection{Thermodynamic Processes}
% ============================================================

\paragraph{Quasi-static processes and the path dependence of heat and work.}
A \textcolor{blue}{quasi-static process} passes through equilibrium states
only, so the state variables are defined along the whole path. Heat and work
then follow from integrals along the path.
% FROM: delta Q and delta W are path-dependent; delta Q <= T dS (Section 3).

\paragraph{Isothermal, isobaric, isochoric and adiabatic processes.}
Each elementary process fixes one variable. The heat and the work of each
process follow from the response functions.
% TODO: table of the four processes with delta Q, delta W and the response
%       function involved; adiabat versus isotherm in the P--V plane.
\begin{equation}
    \boxed{\left(\frac{\partial P}{\partial V}\right)_{\!S}
    =\frac{C_P}{C_V}\left(\frac{\partial P}{\partial V}\right)_{\!T}.}
\end{equation}

\paragraph{Free expansion.}
In a \textcolor{blue}{free expansion}, a gas expands into vacuum inside an
insulated container, so neither heat nor work is exchanged and $U$ stays
constant. The temperature change measures the deviation from ideal
behavior.
% TODO: derive with the second Maxwell relation; irreversible, so Delta S > 0.
\begin{equation}
    \boxed{\left(\frac{\partial T}{\partial V}\right)_{\!U}
    =-\frac{1}{C_V}\left[T\left(\frac{\partial P}{\partial T}\right)_{\!V}-P\right].}
\end{equation}

\paragraph{Throttling.}
In a \textcolor{blue}{throttling process}, a gas flows steadily through a
porous plug from high to low pressure. The enthalpy stays constant, and the
\textcolor{blue}{Joule--Thomson coefficient} gives the temperature change.
% TODO: show H = const from the work balance on both sides of the plug;
%       derive with the fourth Maxwell relation.
\begin{equation}
    \boxed{\mu_{\text{JT}}\equiv\left(\frac{\partial T}{\partial P}\right)_{\!H}
    =\frac{V}{C_P}\left(T\alpha_P-1\right).}
\end{equation}

\paragraph{Heat engines and the Carnot bound.}
A cyclic process returns to the initial state, so the net work equals the
net heat. The second law bounds the fraction of the absorbed heat that a
cycle converts into work.
% OPTIONAL: Carnot cycle; efficiency bound from the Clausius inequality.
\begin{equation}
    \boxed{\eta\leq1-\frac{T_{\text{low}}}{T_{\text{high}}}.}
\end{equation}

% ============================================================
\subsection{The Ideal Gas}
% ============================================================

\paragraph{Classical limit of the partition function.}
The sum over microstates becomes an integral over phase space for classical
particles. Two assumptions enter. Each microstate occupies a phase-space
volume $h^{3N}$, and the factor $1/N!$ removes the permutations of identical
particles.
% FROM: 1/N! and extensivity (Section 5, Gibbs paradox).
% TODO: state the h^{3N} cell as an assumption, justified later by quantum
%       statistics.
\begin{equation}
    \boxed{Z=\frac{1}{N!\,h^{3N}}\int d^{3N}q\,d^{3N}p\;e^{-\beta\mathcal{H}(q,p)}.}
\end{equation}

\paragraph{Partition function of non-interacting particles.}
The \textcolor{blue}{ideal gas} consists of point particles without
interactions, so the Hamiltonian contains the kinetic energy only. The
momentum integrals are Gaussian, and the position integrals give the
volume. \textcolor{red}{Here, $\lambda\equiv h/\sqrt{2\pi mk_BT}$ denotes the
thermal wavelength.}
\begin{equation}
    \boxed{Z=\frac{V^N}{N!\,\lambda^{3N}}.}
\end{equation}

\paragraph{Equation of state, energy and heat capacity.}
The free energy $F=-k_BT\ln Z$ gives the pressure, and
$U=-\partial\ln Z/\partial\beta$ gives the energy.
% FROM: F = -k_B T ln Z, P = -dF/dV (Section 4).
\begin{equation}
    \boxed{PV=Nk_BT,\qquad U=\tfrac{3}{2}Nk_BT,\qquad C_V=\tfrac{3}{2}Nk_B.}
\end{equation}

\paragraph{Entropy and chemical potential.}
The entropy of the ideal gas is the \textcolor{blue}{Sackur--Tetrode
equation}. The entropy is extensive only with the factor $1/N!$.
\begin{equation}
    \boxed{S=Nk_B\left[\ln\frac{V}{N\lambda^3}+\frac{5}{2}\right],\qquad
    \mu=k_BT\ln\!\left(n\lambda^3\right).}
\end{equation}

\paragraph{Response functions and processes of the ideal gas.}
The general results of this section take simple forms for the ideal gas.
Free expansion and throttling leave the temperature of an ideal gas
unchanged, because the particles do not interact.
% TODO: kappa_S, adiabat P V^gamma = const with gamma = 5/3; Var(N)/N = 1.
\begin{equation}
    \boxed{\alpha_P=\frac{1}{T},\qquad\kappa_T=\frac{1}{P},\qquad
    C_P-C_V=Nk_B,\qquad\mu_{\text{JT}}=0.}
\end{equation}

\paragraph{Validity of the ideal gas.}
The ideal gas requires two conditions. The gas must be dilute, so that the
interactions are negligible, and the thermal wavelength must be small
compared with the particle spacing, $n\lambda^3\ll1$, so that the classical
limit holds.

% ============================================================
\subsection{Interacting Gases}
% ============================================================

\paragraph{The configurational integral.}
Pair interactions $\phi(r)$ leave the momentum integrals unchanged. The
partition function then splits into the ideal-gas part and a
\textcolor{blue}{configurational integral} $Q_N$.
\begin{equation}
    \boxed{Z=\frac{Q_N}{N!\,\lambda^{3N}},\qquad
    Q_N=\int d^{3N}r\;\exp\!\Big[-\beta\sum_{i<j}\phi(r_{ij})\Big].}
\end{equation}

\paragraph{The virial expansion.}
At low density, the equation of state is a power series in $1/v$. The
\textcolor{blue}{second virial coefficient} $B_2(T)$ follows exactly from
the pair potential through the \textcolor{blue}{Mayer function}
$f(r)=e^{-\beta\phi(r)}-1$. This expansion is the rigorous link between the
interactions and the deviation from ideal behavior.
% TODO: derive B_2 from Q_N to first order in the density.
\begin{equation}
    \boxed{\frac{Pv}{k_BT}=1+\frac{B_2(T)}{v}+\cdots,\qquad
    B_2(T)=-\frac{1}{2}\int d^3r\,f(r).}
\end{equation}

\paragraph{The van der Waals equation.}
Two assumptions turn $Q_N$ into a closed form. A hard core excludes the
volume $b$ per particle, and the attraction acts on each particle as a
uniform mean field proportional to the density, with strength $a$.
% TODO: Q_N ~ (V - Nb)^N exp(beta a N^2 / V); compare B_2 = b - a/(k_B T)
%       with the virial result.
\begin{equation}
    \boxed{P=\frac{k_BT}{v-b}-\frac{a}{v^2}.}
\end{equation}

\paragraph{Critical point and corresponding states.}
Below a critical temperature, each van der Waals isotherm has a loop, and
the middle part of the loop has $\kappa_T<0$, which violates stability. The loop shrinks to a horizontal
inflection at the critical point, and the reduced variables
$P_r=P/P_c$, $v_r=v/v_c$ and $T_r=T/T_c$ remove the material constants.
\begin{equation}
    \boxed{v_c=3b,\quad k_BT_c=\frac{8a}{27b},\quad P_c=\frac{a}{27b^2},
    \qquad
    \left(P_r+\frac{3}{v_r^2}\right)\left(3v_r-1\right)=8T_r.}
\end{equation}
% TODO: law of corresponding states; the loop is resolved by phase
%       coexistence in the phase-transition sections (Maxwell construction).

\paragraph{The Dieterici equation.}
The \textcolor{blue}{Dieterici equation} replaces the mean-field attraction
by a Boltzmann factor and has the same second virial coefficient as the
van der Waals equation. The critical compressibility factor
$Z_c=P_cv_c/(k_BT_c)$ tests both models, because $Z_c$ is the same for all
substances in either model. Table~\ref{tab:Zc} compares the models with
experiment.
\begin{equation}
    \boxed{P=\frac{k_BT}{v-b}\exp\!\left(-\frac{a}{k_BTv}\right).}
\end{equation}
\begin{table}[h]
\centering
\caption{Critical compressibility factor $Z_c$ of two model equations and
four fluids.}
\label{tab:Zc}
\begin{tabular}{lc}
    Equation or fluid & $Z_c$ \\ \hline
    van der Waals & $3/8=0.375$ \\
    Dieterici & $2/e^2\approx0.271$ \\
    Argon & 0.291 \\
    Nitrogen & 0.289 \\
    Carbon dioxide & 0.274 \\
    Water & 0.229 \\
\end{tabular}
\end{table}

\paragraph{Joule--Thomson inversion.}
The attraction cools a throttled gas, and the hard core heats the gas. At
low density, the van der Waals gas cools on throttling below the
\textcolor{blue}{inversion temperature}.
% TODO: derive mu_JT ~ (2a/(k_B T) - b)/c_P from B_2.
\begin{equation}
    \boxed{k_BT_{\text{inv}}=\frac{2a}{b}=\frac{27}{4}\,k_BT_c.}
\end{equation}
